\subsection{Characterization of polynomials with zero differential}

PROPOSITION 6. --- Let K be a commutative ring, A the polynomial algebra \(K[X_{i}]_{i \in I}\) and S the set of elements F of A such that dF = 0.

\begin{enumerate}
\def\labelenumi{\alph{enumi})}
\item
  If \(\mathbf{K}\) is a ring of characteristic 0, then \(\mathbf{S} = \mathbf{K}\) .
\item
  If \(\mathbf{K}\) is a ring of characteristic \(p \neq 0\) , then \(\mathbf{S} = \mathbf{K}[\mathbf{X}_i^p]_{i \in I}\) ; if moreover \(\mathbf{K}\) is perfect, then \(\mathbf{K} = \mathbf{A}^p\) .
\end{enumerate}

The mapping \(F \mapsto dF\) of A into the module \(\Omega_{\mathbf{K}}(\mathbf{A})\) of K-differentials of A is K-linear and satisfies the relation

\[
d \left(\mathrm{FF} ^ {\prime}\right) = \mathrm{F} \cdot d \mathrm{F} ^ {\prime} + \mathrm{F} ^ {\prime} \cdot d \mathrm{F}
\]

(III, p.~569). Therefore S is a subalgebra.

When K is of characteristic \(p \neq 0\) , put \(T = K[X_{i}^{p}]_{i \in I}\) ; we thus have \(T = A^{p}\) if K is perfect (V, p.~4, Prop. 2); moreover, we have \(d(X_{i}^{p}) = pX_{i}^{p-1}\) . \(dX_{i} = 0\) for all \(i \in I\) , hence the subalgebra S of A contains T. If K is of characteristic 0, we put T = K, and still find that \(T \subset S\) . It remains to show that S is contained in T.

For every finite subset \(J\) of \(I\) let \(A_{J}\) be the subalgebra of \(A\) generated by the family \((X_{j})_{j\in J}\) . We have \(A_{\emptyset} = K\) and \(A = \bigcup_{J\in I}A_{J}\) ; so it suffices to prove the

relation \(S \cap A_{J} \subset T\) , which we shall accomplish by induction on the cardinal of J. Thus let J be a finite subset of I such that \(S \cap A_{J} \subset T\) , let i be an element of I - J and \(J' = J \cup \{i\}\) . Every element F of \(A_{J}\) may be written in just one way in the form

\[
\mathrm{F} = \sum_ {n = 0} ^ {\infty} \mathrm{F} _ {n} \cdot \mathrm{X} _ {i} ^ {n},\tag{5}
\]

with \(\mathbf{F}_n\in \mathbf{A}_{\mathrm{J}}\) for all \(n\geqslant 0\) , and then

\[
d \mathrm{F} = \sum_ {n = 0} ^ {\infty} \mathrm{X} _ {i} ^ {n} \cdot d \mathrm{F} _ {n} + \sum_ {n = 0} ^ {\infty} n \mathrm{X} _ {i} ^ {n - 1} \mathrm{F} _ {n} \cdot d \mathrm{X} _ {i}.\tag{6}
\]

Suppose that F belongs to S; the family \((d\mathbf{X}_{r})_{r\in I}\) is a basis of the A-module \(\Omega_{\mathrm{K}}(\mathbf{A})\) (III, p.~570) and \(dF_{n}\) is a linear combination of the differentials \(dX_{j}\) for \(j\in J\) because \(F_{n}\in A_{J}=K[X_{j}]_{j\in J}\) . By (6) we then have \(dF_{n}=0\) and \(nF_{n}=0\) for each integer \(n\geqslant0\) . By the induction hypothesis \(F_{n}\in T\) for all n, since \(dF_{n}=0\) .

\begin{enumerate}
\def\labelenumi{\alph{enumi})}
\item
  If K is of characteristic 0, we have \(nF_{n} = 0\) for all \(n \geqslant 1\) , whence \(F_{n} = 0\) by Prop. 1 (V, p.~2); so we have \(F = F_{0}\) , whence \(F \in T\) .
\item
  If \(\mathbf{K}\) is of characteristic \(p \neq 0\) , then \(\mathbf{A}\) is an algebra over the field \(\mathbf{F}_p\) and the relation \(n\mathrm{F}_n = 0\) implies \(\mathbf{F}_n = 0\) for every integer \(n\) not divisible by \(p\) . So we have \(\mathbf{F} = \sum_{m=0}^{\infty} \mathbf{F}_{mp} \mathbf{X}_i^{mp}\) , whence \(\mathbf{F} \in \mathbf{T}\) .
\end{enumerate}

Remark. --- We still have S = K when the additive group of K is torsion-free; this follows from the above proof or from Remark 5 of V, p.~3.

COROLLARY. --- Let K be a field and F(X) a polynomial with coefficients in K, whose derivative F'(X) is zero.

\begin{enumerate}
\def\labelenumi{\alph{enumi})}
\item
  If \(\mathbf{K}\) is of characteristic 0, then \(\mathbf{F} \in \mathbf{K}\) .
\item
  If \(\mathbf{K}\) is of characteristic \(p \neq 0\) , there exists a polynomial \(\mathbf{G}(\mathbf{X})\) such that \(\mathbf{F}(\mathbf{X}) = \mathbf{G}(\mathbf{X}^p)\) .
\end{enumerate}

For we have \(d\mathbf{F} = \mathbf{F}' \cdot d\mathbf{X} = 0\) .

